\documentclass[11pt,a4paper]{article}
\usepackage[utf8]{inputenc}
\usepackage[spanish,es-tabla]{babel}
\usepackage{amsmath,amsfonts,amssymb}
\usepackage{geometry}
\usepackage{graphicx}
\usepackage{xcolor}

\usepackage{enumitem}

\geometry{margin=2.5cm}

\title{\textbf{Soluci\'on Examen 2019-2020\\TAF ISC - UEC 2\\18/11/2019}}
\author{}
\date{}

\begin{document}

\maketitle

\section*{Ejercicio 1}

Consideramos una transmisi\'on sobre frecuencia porteuse con se\~nal emitida:
\[
s(t) = \sum_{k\in\mathbb{Z}} a_k h(t-kT_s)\cos(2\pi\nu_0 t + \theta_0) - b_k h(t-kT_s)\sin(2\pi\nu_0 t + \theta_0)
\]

\subsection*{1.1 Preguntas de curso}

\subsection*{Pregunta 1: Criterio TSWOC}

\textbf{Enunciado:} Enunciar el criterio Time-Shifted-Waveform Orthogonality Criterion (TSWOC) en funci\'on de $h(t)$.

\vspace{0.3cm}
\textbf{Soluci\'on:}

El criterio TSWOC establece que:
\[
\boxed{\int_{\mathbb{R}} h(t) h(t - kT_s)\, dt = \delta_{k,0}}
\]

Donde:
\begin{itemize}
    \item $h(t)$ es la respuesta impulsional del filtro de conformaci\'on
    \item $T_s$ es el periodo s\'imbolo
    \item $\delta_{k,0}$ es la delta de Kronecker: vale 1 si $k=0$ y 0 en caso contrario
\end{itemize}

Esto significa que las versiones desplazadas del pulso deben ser ortogonales entre s\'i.



\subsection*{Pregunta 2: Consecuencias del TSWOC}

\textbf{Enunciado:} ¿Qu\'e garantiza este criterio para una transmisi\'on sobre un canal gaussiano bajo la hip\'otesis de un receptor \'optimo? (3 respuestas)

\vspace{0.3cm}
\textbf{Soluci\'on:}

El criterio TSWOC garantiza tres propiedades fundamentales:

\begin{enumerate}
    \item \textbf{Ausencia de interferencia entre s\'imbolos (ISI):} Los s\'imbolos transmitidos no interfieren entre s\'i en los instantes de muestreo.
    
    \item \textbf{Detecci\'on s\'imbolo por s\'imbolo \'optima:} Se puede tomar la decisi\'on sobre cada s\'imbolo de forma independiente.
    
    \item \textbf{Maximizaci\'on del SNR:} Se maximiza la relaci\'on se\~nal a ruido a la salida del filtro adaptado.
\end{enumerate}



\subsection*{1.2 C\'alculos de par\'ametros de la transmisi\'on}

\textbf{Datos del problema:}
\begin{itemize}
    \item Modulaci\'on: MAQ-256 (256-QAM)
    \item Rapidez de modulaci\'on: $D_s = 256$ kbauds = $256 \times 10^3$ s\'imbolos/s
    \item Frecuencia portadora: $\nu_0 = 2$ GHz = $2 \times 10^9$ Hz
    \item Filtro: Raised Cosine (cosinus sur\'elev\'e) con coeficiente de roll-off $\beta = \frac{3}{4}$
\end{itemize}

\subsection*{Pregunta 1: C\'alculo del d\'ebit binaire}

\textbf{Enunciado:} Calcular el d\'ebito binario.

\vspace{0.3cm}
\textbf{Soluci\'on:}

El d\'ebito binario se relaciona con la rapidez de modulaci\'on mediante:
\[
D_b = D_s \times \log_2(M)
\]

Donde $M$ es el orden de la modulaci\'on (n\'umero de s\'imbolos en el alfabeto).

Para MAQ-256: $M = 256 = 2^8$

Por lo tanto:
\[
D_b = 256 \times 10^3 \times 8 = \boxed{2.048 \text{ Mbps}}
\]

\textbf{Interpretaci\'on:} Cada s\'imbolo transporta 8 bits de informaci\'on.



\subsection*{Pregunta 2: C\'alculo de la banda ocupada}

\textbf{Enunciado:} Calcular la anchura de banda ocupada.

\vspace{0.3cm}
\textbf{Soluci\'on:}

Para un filtro raised cosine con roll-off $\beta$, la banda ocupada es:
\[
B_{occ} = D_s \times (1 + \beta)
\]

Sustituyendo valores:
\[
B_{occ} = 256 \times 10^3 \times \left(1 + \frac{3}{4}\right) = 256 \times 10^3 \times 1.75
\]
\[
B_{occ} = \boxed{448 \text{ kHz}}
\]

\textbf{Nota:} La banda m\'inima te\'orica ser\'ia $D_s = 256$ kHz, pero el exceso de banda $\beta$ la incrementa.



\subsection*{Pregunta 3: C\'alculo del exceso de banda}

\textbf{Enunciado:} Calcular el exceso de banda.

\vspace{0.3cm}
\textbf{Soluci\'on:}

El exceso de banda se define como:
\[
\Delta B = B_{occ} - B_{min} = B_{occ} - D_s
\]

O equivalentemente:
\[
\Delta B = D_s \times \beta
\]

Sustituyendo:
\[
\Delta B = 256 \times 10^3 \times \frac{3}{4} = \boxed{192 \text{ kHz}}
\]

\textbf{Interpretaci\'on:} El filtro raised cosine utiliza un 75\% m\'as de banda que el m\'inimo te\'orico, pero a cambio proporciona mejor comportamiento en presencia de errores de sincronizaci\'on.



\subsection*{Pregunta 4: C\'alculo de la eficiencia espectral}

\textbf{Enunciado:} Calcular la eficacia espectral de la transmisi\'on.

\vspace{0.3cm}
\textbf{Soluci\'on:}

La eficiencia espectral se define como:
\[
\eta = \frac{D_b}{B_{occ}} \text{ [bits/s/Hz]}
\]

Sustituyendo los valores calculados:
\[
\eta = \frac{2.048 \times 10^6}{448 \times 10^3} = \frac{2048}{448} = \boxed{4.57 \text{ bits/s/Hz}}
\]

\textbf{An\'alisis:}
\begin{itemize}
    \item Valor te\'orico m\'aximo para esta modulaci\'on: $\frac{\log_2(M)}{1+\beta} = \frac{8}{1.75} = 4.57$ bits/s/Hz ✓
    \item L\'imite de Shannon para canal ideal: $\eta_{Shannon} \approx 10-12$ bits/s/Hz (con SNR alto)
    \item Valor pr\'actico: moderado, compromiso entre eficiencia y robustez
\end{itemize}



\newpage
\section*{Problema 2}

Consideramos un alfabeto de dos estados $\mathcal{Q} = \{q_1, q_2\}$ definido por:
\begin{align*}
q_1(t) &= Ah(t)\cos(2\pi\nu_0 t + \theta_0) \\
q_2(t) &= A\cos(\Delta\theta)h(t)\cos(2\pi\nu_0 t + \theta_0) - A\sin(\Delta\theta)h(t)\sin(2\pi\nu_0 t + \theta_0)
\end{align*}

con $0 < \Delta\theta \leq \pi$.

\textbf{Hip\'otesis:}
\begin{itemize}
    \item Banda estrecha: el ancho de banda de $h(t)$ es muy peque\~no comparado con $\nu_0$
    \item Modulaci\'on de un bit $m_0$ con $\Pr(m_0 = 0) = \Pr(m_0 = 1) = \frac{1}{2}$
    \item Canal AWGN: $r(t) = q(t) + n(t)$ con $n(t) \sim \mathcal{N}(0, \frac{N_0}{2})$
\end{itemize}

\subsection*{Pregunta 1: Aplicaci\'on del criterio MAP}

\textbf{Enunciado:} Aplicar el criterio de detecci\'on MAP y dar la regla de decisi\'on.

Criterio MAP:
\[
\hat{q} = \arg\max_{q_j \in \mathcal{Q}} \left\langle r(t), q_j(t) \right\rangle + N_0 \Pr(q = q_j) - \frac{E_j}{2}
\]

\vspace{0.3cm}
\textbf{Soluci\'on:}

\textbf{Paso 1:} Calcular las energ\'ias de los s\'imbolos.

Asumiendo $\int_{\mathbb{R}} |h(t)|^2\, dt = 1$ (normalizaci\'on):

\[
E_1 = \int_{\mathbb{R}} |q_1(t)|^2\, dt = A^2 \int_{\mathbb{R}} h^2(t)\cos^2(2\pi\nu_0 t + \theta_0)\, dt \approx \frac{A^2}{2}
\]

Para $q_2(t)$, usando notaci\'on compleja:
\[
q_2(t) = A\cos(\Delta\theta)h(t)\cos(2\pi\nu_0 t + \theta_0) - A\sin(\Delta\theta)h(t)\sin(2\pi\nu_0 t + \theta_0)
\]

\[
E_2 = A^2[\cos^2(\Delta\theta) + \sin^2(\Delta\theta)] \int_{\mathbb{R}} h^2(t) \cdot \frac{1}{2}\, dt = \frac{A^2}{2}
\]

Por lo tanto: $E_1 = E_2 = \frac{A^2}{2}$

\textbf{Paso 2:} Como las probabilidades son iguales y las energ\'ias tambi\'en, el criterio se simplifica:
\[
\hat{q} = \arg\max_{q_j \in \{q_1, q_2\}} \left\langle r(t), q_j(t) \right\rangle
\]

\textbf{Paso 3:} Calcular los productos escalares con las se\~nales propuestas:

Definimos:
\begin{align*}
z_c &= \left\langle r(t), Ah(t)\cos(2\pi\nu_0 t + \theta_0) \right\rangle \\
z_s &= -\left\langle r(t), Ah(t)\sin(2\pi\nu_0 t + \theta_0) \right\rangle
\end{align*}

Entonces:
\begin{align*}
\left\langle r(t), q_1(t) \right\rangle &= z_c \\
\left\langle r(t), q_2(t) \right\rangle &= z_c \cos(\Delta\theta) + z_s \sin(\Delta\theta)
\end{align*}

\textbf{Regla de decisi\'on:}
\[
\boxed{
\begin{cases}
\hat{m}_0 = 0 & \text{si } z_c > z_c \cos(\Delta\theta) + z_s \sin(\Delta\theta) \\
\hat{m}_0 = 1 & \text{en caso contrario}
\end{cases}
}
\]

Simplificando:
\[
\boxed{
\begin{cases}
\hat{m}_0 = 0 & \text{si } z_c(1 - \cos(\Delta\theta)) > z_s \sin(\Delta\theta) \\
\hat{m}_0 = 1 & \text{en caso contrario}
\end{cases}
}
\]



\subsection*{Pregunta 2: Estructura del receptor \'optimo}

\textbf{Enunciado:} Representar la estructura del receptor \'optimo asociada a esta regla.

\vspace{0.3cm}
\textbf{Soluci\'on:}

El receptor \'optimo consiste en:

\begin{center}
\begin{tabular}{c}
\texttt{r(t)} \\
$\downarrow$ \\
\fbox{$\times \cos(2\pi\nu_0 t + \theta_0)$} $\rightarrow$ \fbox{Filtro $h(t_0 - t)$} $\rightarrow$ \fbox{Muestreo $t_0$} $\rightarrow z_c$ \\
$\downarrow$ \\
\fbox{$\times (-\sin(2\pi\nu_0 t + \theta_0))$} $\rightarrow$ \fbox{Filtro $h(t_0 - t)$} $\rightarrow$ \fbox{Muestreo $t_0$} $\rightarrow z_s$ \\
\end{tabular}
\end{center}

\vspace{0.3cm}
\begin{center}
$z_c, z_s$ $\rightarrow$ \fbox{Dispositivo de decisi\'on} $\rightarrow \hat{m}_0$
\end{center}

\textbf{Componentes:}
\begin{itemize}
    \item \textbf{Demodulador I (In-phase):} Multiplica por $\cos(2\pi\nu_0 t + \theta_0)$ para extraer componente en fase
    \item \textbf{Demodulador Q (Quadrature):} Multiplica por $-\sin(2\pi\nu_0 t + \theta_0)$ para extraer componente en cuadratura
    \item \textbf{Filtros adaptados:} $h(t_0 - t)$ maximizan el SNR
    \item \textbf{Muestreadores:} En instante \'optimo $t_0$
    \item \textbf{Decisor:} Aplica la regla $z_c(1-\cos\Delta\theta) \stackrel{?}{>} z_s\sin\Delta\theta$
\end{itemize}



\subsection*{Pregunta 3: Constelaci\'on recibida}

\textbf{Enunciado:} Representar la constelaci\'on recibida para un valor arbitrario de $\Delta\theta$.

\vspace{0.3cm}
\textbf{Soluci\'on:}

En el plano $(z_c, z_s)$, las se\~nales recibidas (sin ruido) son:

Para $m_0 = 0$ (se transmite $q_1$):
\[
z_c = \frac{A^2}{2}, \quad z_s = 0
\]

Para $m_0 = 1$ (se transmite $q_2$):
\[
z_c = \frac{A^2}{2}\cos(\Delta\theta), \quad z_s = \frac{A^2}{2}\sin(\Delta\theta)
\]

\textbf{Representaci\'on gr\'afica:}

\begin{center}
\begin{tikzpicture}[scale=2]
% Ejes
\draw[->] (-0.5,0) -- (2,0) node[right] {$z_c$};
\draw[->] (0,-0.3) -- (0,1.5) node[above] {$z_s$};

% \'Angulo arbitrario Delta theta = 60 grados
\def\angle{60}
\def\amplitude{1.5}

% Punto q1
\fill[blue] (\amplitude,0) circle (2pt) node[below right] {$q_1$ ($m_0=0$)};

% Punto q2
\fill[red] ({\amplitude*cos(\angle)},{\amplitude*sin(\angle)}) circle (2pt) node[above right] {$q_2$ ($m_0=1$)};

% L\'ineas auxiliares
\draw[dashed, gray] (0,0) -- (\amplitude,0);
\draw[dashed, gray] (0,0) -- ({\amplitude*cos(\angle)},{\amplitude*sin(\angle)});

% \'Angulo
\draw[<->, thick] (0.5,0) arc (0:\angle:0.5) node[midway, right] {$\Delta\theta$};

% Frontera de decisi\'on
\draw[thick, green!70!black] (0.8,-0.2) -- (1.2,1.2) node[pos=0.8, right] {Frontera};

\end{tikzpicture}
\end{center}

\textbf{Observaciones:}
\begin{itemize}
    \item Los dos s\'imbolos tienen la misma energ\'ia: $|q_1| = |q_2| = \frac{A^2}{2}$
    \item Est\'an separados por un \'angulo $\Delta\theta$
    \item La frontera de decisi\'on es la mediatriz del segmento que une ambos puntos
\end{itemize}



\subsection*{Pregunta 4: Valor \'optimo de $\Delta\theta$}

\textbf{Enunciado:} Deducir el valor \'optimo de $\Delta\theta$ que minimiza la probabilidad de error binaria (justificaci\'on gr\'afica).

\vspace{0.3cm}
\textbf{Soluci\'on:}

\textbf{An\'alisis:}

La probabilidad de error en un canal AWGN depende de la distancia euclidiana entre los dos s\'imbolos:

\[
d_{12}^2 = |q_1 - q_2|^2 = \int_{\mathbb{R}} [q_1(t) - q_2(t)]^2\, dt
\]

Desarrollando:
\begin{align*}
d_{12}^2 &= E_1 + E_2 - 2\langle q_1, q_2 \rangle \\
&= \frac{A^2}{2} + \frac{A^2}{2} - 2 \cdot \frac{A^2}{2}\cos(\Delta\theta) \\
&= A^2(1 - \cos(\Delta\theta))
\end{align*}

Para minimizar la probabilidad de error, debemos \textbf{maximizar} la distancia $d_{12}$.

Como $d_{12}^2 = A^2(1 - \cos(\Delta\theta))$, esto se maximiza cuando $\cos(\Delta\theta)$ es m\'inimo.

En el intervalo $(0, \pi]$, el coseno alcanza su m\'inimo en $\Delta\theta = \pi$:

\[
\boxed{\Delta\theta_{opt} = \pi}
\]

\textbf{Justificaci\'on gr\'afica:}

Con $\Delta\theta = \pi$, los dos s\'imbolos son antipodales:
\begin{itemize}
    \item $q_1$: punto en $(+\frac{A^2}{2}, 0)$
    \item $q_2$: punto en $(-\frac{A^2}{2}, 0)$
\end{itemize}

La distancia entre ellos es $d_{12} = A^2$, m\'axima posible.

\textbf{Interpretaci\'on f\'isica:}

$\Delta\theta = \pi$ corresponde a una modulaci\'on BPSK (Binary Phase Shift Keying), que es la modulaci\'on binaria \'optima para canal AWGN.

\textbf{Probabilidad de error resultante:}
\[
P_e = Q\left(\sqrt{\frac{2E_b}{N_0}}\right) = Q\left(\sqrt{\frac{A^2}{N_0}}\right)
\]

donde $Q(\cdot)$ es la funci\'on Q de Gauss.



\vspace{1cm}
\begin{center}
\rule{0.8\linewidth}{0.4pt}\\
\textit{Fin del examen}
\end{center}

\end{document}
